\section{3.3~V converter derivation}\label{sec:3v3-converter}

This section derives the required characteristics of the 3.3~V rail and verifies the
populated Rev-A implementation: U411 (AP63203WU, \cref{sec:ap632xx}) with L411
(\SI{4.7}{\micro\henry}), C411 (\SI{10}{\micro\farad} input), C412 (\SI{100}{\nano\farad}
BST), and C413/C414 ($2\times\SI{22}{\micro\farad}$ output). The exported netlist confirms
the \texttt{+3V3} net contains exactly: the converter output components (L411, C413,
C414), the ESP32-S3-WROOM-1 module's supply pin (U501 pin 2) with its local decoupling
C501/C502, U411's FB pin (tied to the output, as on all fixed variants), and five
resistive sources feeding other nets --- R206 (LM5155 PGOOD pull-up), R501 (ESP32 EN
network), R502/R503 (UART activity-LED feeds), and R411 (rail indicator LED). The
converter is the same IC family, footprint set, and support-component topology as the 5~V
rail of \cref{sec:5v-converter}, and the generic buck method established there ---
selection equations, mode-boundary logic, transient bound, RMS estimates --- is applied
here without re-derivation. What is genuinely different on this rail, and where this
section spends its effort, is threefold: the load is large and transient-heavy (Wi-Fi TX
bursts, not static logic current), the converter's input is a two-diode OR of
$V_{\mathrm{BUS}}$ and an externally supplied UART-header voltage rather than a direct
rail connection, and the rail sources a status pull-up (PGOOD) whose behavior must be
consistent with the rail's own sequencing.

\subsection{Load budget}\label{sec:3v3-load}

The dominant load is the MCU module. From \cref{sec:esp32}: the worst published sustained
condition is Wi-Fi TX at \SI{355}{\milli\ampere} (802.11b, \SI{20.5}{dBm}, rated at
100\,\% TX duty cycle --- a sustained worst case, not an instantaneous spike atop a
sustained figure), and Espressif's own regulator-sizing requirement is an external supply
capable of $\geq\SI{0.5}{\ampere}$. Typical operation sits far lower: modem-sleep spans
\SIrange{13.2}{91.7}{\milli\ampere} depending on CPU state and clock, with a
\SI{10}{\milli\ampere} adder during flash access; light- and deep-sleep are in the
\si{\micro\ampere} range. The remaining loads are checked for negligibility rather than
derived at length:

\begin{itemize}
\item \emph{BCD drive lines.} Sixteen GPIOs drive the four K155ID1s' BCD inputs. A TTL
input driven high draws $I_{IH} \leq \SI{40}{\micro\ampere}$ (A input) or
\SI{80}{\micro\ampere} (B/C/D inputs) from the driving pin (\cref{sec:bcd-driver}); a
line driven low \emph{sinks} current sourced from the drivers' 5~V rail and adds nothing
to VDD33. Worst case, all sixteen lines high:
\begin{align*}
  I_{\mathrm{BCD}} &= 4 \times (40 + 3\times80)\,\si{\micro\ampere}
    = \SI{1.12}{\milli\ampere},
\end{align*}
0.3\,\% of the Wi-Fi figure. Negligible, as expected.
\item \emph{Indicator LEDs.} R411 (\SI{1.5}{\kilo\ohm}) feeds D413, the rail's power
indicator, continuously; R502/R503 (\SI{1.5}{\kilo\ohm} each) feed D501/D502, activity
LEDs whose cathodes sit on the RXD0/TXD0 nets and which conduct only while the
corresponding UART line is driven low (the UART idle state is high, so idle draw is
zero). \emph{Assumption:} LED forward voltage $\approx\SI{2.0}{\volt}$ (generic 0805
part, color unrecorded). Then
\begin{align*}
  I_{\mathrm{D413}} &= \frac{3.3-2.0}{\SI{1.5}{\kilo\ohm}} = \SI{0.87}{\milli\ampere}
    \quad\text{(continuous)}, \\
  I_{\mathrm{D501,D502}} &\approx \frac{3.3-2.0-0.1}{\SI{1.5}{\kilo\ohm}}
    = \SI{0.8}{\milli\ampere} \quad\text{(per line, while low).}
\end{align*}
(D413's value field reads ``RX,'' evidently copied from the activity LEDs; electrically
it is the rail indicator --- a cosmetic schematic-field inconsistency, flagged for BOM
hygiene only.)
\item \emph{Pull-ups.} R206 (PGOOD) contributes at most \SI{33}{\micro\ampere} and only
during power-not-good windows; R501 (EN) conducts only transiently. Both are quantified
in \cref{sec:3v3-house}.
\end{itemize}

The design totals are therefore
\begin{align*}
  I_{\mathrm{load,max}} &= 355 + 1.12 + 0.87 + 2\times0.8 + 0.03
    \approx \SI{359}{\milli\ampere} \approx \SI{0.36}{\ampere},\\
  P_{\mathrm{out,max}} &= 3.3 \times 0.359 = \SI{1.19}{\watt},
\end{align*}
the sustained worst case, with Espressif's $\geq\SI{0.5}{\ampere}$ capability requirement adopted as the
converter-sizing floor (margin checks below are run at both \SI{0.36}{\ampere} and
\SI{0.5}{\ampere}). Against the AP63203's \SI{2}{\ampere} rating the capability
requirement is met $4\times$ over. For the UART-powered scenario of
\cref{sec:3v3-diodeor} --- an MCU in download/console mode, radio off --- the applicable
draw is the modem-sleep band plus flash access and LEDs: roughly
\SIrange{30}{100}{\milli\ampere}, with \SI{50}{\milli\ampere} carried as the nominal
estimate. One open item is inherited from \cref{sec:esp32} and not resolvable here:
if the installed module SKU has PSRAM, active-mode draw may exceed the tabulated
\SI{355}{\milli\ampere} by an amount Espressif does not quantify. The $4\times$ headroom
between the \SI{0.5}{\ampere} requirement and the \SI{2}{\ampere} part absorbs any
plausible adder, but the published figures remain the design point.

\subsection{Operating point and conduction mode}\label{sec:3v3-op}

Unlike U401, whose IN/EN tie directly to $V_{\mathrm{BUS}}$, U411's IN and EN share a
common node fed through Schottky diode D411 from $V_{\mathrm{BUS}}$ (the OR topology
itself is treated in \cref{sec:3v3-diodeor}; this subsection covers the normal,
$V_{\mathrm{BUS}}$-fed case). \emph{Assumption} (generic SOD-123 Schottky, no part
number recorded): $V_F \approx \SI{0.3}{\volt}$ at \SI{50}{\milli\ampere},
\SI{0.4}{\volt} at \SIrange{200}{300}{\milli\ampere}, \SI{0.45}{\volt} at
\SI{0.5}{\ampere}; reverse leakage $\leq\SI{100}{\micro\ampere}$ warm at \SI{13}{\volt}.
The converter's input node then spans
\begin{align*}
  V_N^{\min} &= 8.5 - 0.4 = \SI{8.1}{\volt} \quad\text{(full load, minimum bus)}, \\
  V_N^{\max} &= 13 - 0.2 \approx \SI{12.8}{\volt} \quad\text{(light load, maximum bus)},
\end{align*}
inside the \SIrange{3.8}{32}{\volt} recommended range with the same style of margin as
the 5~V rail; EN\,=\,IN is a sanctioned always-on configuration, and both pins' absolute
maximum is \SI{35}{\volt} (\cref{sec:ap632xx}) --- a $2.7\times$ margin over
$V_N^{\max}$. The ideal duty cycle spans
\begin{align*}
  D_{\min} &= \frac{3.3}{12.8} = 0.258, &
  D_{\max} &= \frac{3.3}{8.1} = 0.407,
\end{align*}
never crossing 0.5, so even the internal slope compensation is unstressed. Conduction
corrections ($I R_{DS(\mathrm{on})} \approx \SI{45}{\milli\volt}$ at full load, 1.4\,\%
of $V_{OUT}$) are ignored as in \cref{sec:5v-op}.

\paragraph{Conduction mode.} With the populated \SI{4.7}{\micro\henry} at
$f_{SW}=\SI{1.1}{\mega\hertz}$, the fixed-frequency ripple current is
\begin{align*}
  \Delta I_L \big|_{\SI{12.8}{\volt}} &=
    \frac{3.3\,(12.8-3.3)}{12.8 \times \SI{4.7}{\micro\henry} \times \SI{1.1}{\mega\hertz}}
    = \frac{31.35}{66.2} = \SI{0.474}{\ampere}, \\
  \Delta I_L \big|_{\SI{8.1}{\volt}} &=
    \frac{3.3\,(8.1-3.3)}{8.1 \times \SI{4.7}{\micro\henry} \times \SI{1.1}{\mega\hertz}}
    = \frac{15.84}{41.9} = \SI{0.378}{\ampere},
\end{align*}
so the CCM boundary ($\Delta I_L/2$) is \SIrange{189}{237}{\milli\ampere} across the
input window. Where the 5~V rail's load sits permanently below its boundary, this rail
\emph{crosses} it in operation: modem-sleep currents (\SIrange{13}{92}{\milli\ampere})
sit in PFM, and a Wi-Fi TX burst at \SI{355}{\milli\ampere} sits in CCM. The part
manages the PFM/PWM transition automatically and continuously --- the mechanism, its
\SI{450}{\milli\ampere} PFM peak-current clamp, and the zero-cross blocking are as
described in \cref{sec:5v-op} --- so the crossing imposes no component choice; its one
real consequence, transient droop at the mode transition, is bounded in
\cref{sec:3v3-cout}. The datasheet's typical-efficiency curve for
$V_{IN}=\SI{12}{\volt}$, $V_{OUT}=\SI{3.3}{\volt}$ reads approximately 87--90\,\% over
the \SIrange{50}{500}{\milli\ampere} span (\ds{AP632XX.pdf}{2}); $\eta=0.85$ is carried
as the conservative figure, matching \cref{sec:5v-op}. Average input current at full
load is then
\begin{align*}
  I_{IN} &= \frac{P_{\mathrm{out,max}}}{\eta\,V_N^{\min}}
          = \frac{1.19}{0.85 \times 8.1} = \SI{0.17}{\ampere}
  \quad (\SI{0.24}{\ampere}\ \text{at the \SI{0.5}{\ampere} capability point}),
\end{align*}
drawn from $V_{\mathrm{BUS}}$ through D411 --- the largest of the three logic-rail
loads on the bus budget, and the diode adds $I_{IN} V_F \approx \SI{70}{\milli\watt}$
(\SI{96}{\milli\watt} at the capability point) of series loss, about 5\,\% of
throughput at minimum bus voltage.

\subsection{Minimum on-time}\label{sec:3v3-ontime}

The binding corner is $V_N^{\max}=\SI{12.8}{\volt}$. At the $+6\,\%$ spread-spectrum
extreme ($f_{SW}=\SI{1.166}{\mega\hertz}$):
\begin{align*}
  t_{\mathrm{on}} &= \frac{D_{\min}}{f_{SW}}
    = \frac{0.258}{\SI{1.166}{\mega\hertz}} = \SI{221}{\nano\second}
    \quad\text{vs.}\quad t_{\mathrm{on,min}} = \SI{80}{\nano\second}
    \;\;(\cref{sec:ap632xx}),
\end{align*}
a $2.8\times$ margin; the PFM pulse on-time, $L\times\SI{0.45}{\ampere}/(V_N-V_{OUT}) =
\SI{223}{\nano\second}$ at the same corner, is the same figure. The margin is thinner
than the 5~V rail's $3.3$--$4.1\times$ (\cref{sec:5v-ontime}) because a 3.3~V output
yields a smaller duty cycle at the same input, but for the constraint to bind the input
would need $V_N = 3.3/(t_{\mathrm{on,min}} f_{SW}) \approx \SI{35}{\volt}$ --- the
absolute-maximum rating itself, unreachable from a 13~V bus. Match.

\subsection{Inductor}\label{sec:3v3-inductor}

The datasheet selection equation with its prescribed ripple band ($\Delta I_L$ =
30--50\,\% of the \SI{2}{\ampere} rating, i.e.\ \SIrange{0.6}{1.0}{\ampere};
\cref{sec:ap632xx}), at the maximum-ripple corner $V_N=\SI{12.8}{\volt}$:
\begin{align*}
  L\big|_{\Delta I_L = \SI{1.0}{\ampere}} &=
    \frac{3.3\,(12.8-3.3)}{12.8 \times 1.0 \times \SI{1.1}{\mega\hertz}}
    = \SI{2.23}{\micro\henry},
  &
  L\big|_{\Delta I_L = \SI{0.6}{\ampere}} &=
    \frac{31.35}{12.8 \times 0.6 \times \SI{1.1}{\mega\hertz}} = \SI{3.71}{\micro\henry},
\end{align*}
and \SIrange{1.78}{2.96}{\micro\henry} at \SI{8.1}{\volt}. The datasheet's own Table-2
BOM value for the AP63203 is \SI{3.9}{\micro\henry} (\cref{sec:ap632xx}). The populated
value is \SI{4.7}{\micro\henry} --- the \emph{AP63205's} Table-3 value, the same part
and footprint as the 5~V rail's L401. Unlike the 5~V rail, then, this converter is not
the reference BOM applied verbatim: the inductor deviates from the variant-specific
recommendation, evidently to share one line item across both rails. The consequences
all fall in the benign direction. Ripple is \SI{0.474}{\ampere} (23.7\,\% of
\SI{2}{\ampere}) at \SI{12.8}{\volt} and \SI{0.378}{\ampere} (18.9\,\%) at
\SI{8.1}{\volt} --- below the 30--50\,\% band, which lowers the CCM peak, lowers the
output ripple, and improves light-load efficiency (the datasheet's own endorsement of
larger inductance, apt for a clock that idles at modem-sleep currents for essentially
its whole life). The value remains inside the general \SIrange{2.2}{10}{\micro\henry}
recommendation. The cost is a slightly slower current slew into a load step, which the
transient check of \cref{sec:3v3-cout} absorbs with wide margin. Verdict: acceptable ---
a deliberate cross-rail reuse, noted as a deviation from the per-variant BOM table
rather than a mismatch.

\paragraph{Saturation current.} The peak inductor current in normal operation is
\begin{align*}
  I_{\mathrm{pk}} &= I_{\mathrm{load,max}} + \frac{\Delta I_L}{2}
    = 0.359 + 0.237 = \SI{0.60}{\ampere}
  \quad (\SI{0.74}{\ampere}\ \text{at \SI{0.5}{\ampere}}),
\end{align*}
against the datasheet's $\geq$35\,\% rule: $1.35 \times 0.359 = \SI{0.48}{\ampere}$
($1.35 \times 0.5 = \SI{0.68}{\ampere}$ at the capability point). Under the same
stated assumption as \cref{sec:5v-inductor} --- $I_{\mathrm{sat}} \geq \SI{3.0}{\ampere}$,
DCR $\leq\SI{50}{\milli\ohm}$ for a commodity 6045-class \SI{4.7}{\micro\henry} shielded
part --- the rating exceeds the worst normal peak by $\approx 5\times$ (against a real,
sub-ampere operating peak this time, not a PFM clamp), and covers the
\SI{3.1}{\ampere} maximum fault peak at or near hard saturation with hiccup protection
bounding dwell, exactly as argued there. Match under the stated assumption; the actual
part's rating is a confirm-at-BOM item shared with L401.

\subsection{Current-limit margin}\label{sec:3v3-ilim}

Against the \SI{2.5}{\ampere} minimum HS peak limit (\cref{sec:ap632xx}):
\begin{align*}
  \frac{I_{\mathrm{lim}}^{\min}}{I_{\mathrm{pk}}} &= \frac{2.5}{0.60} = 4.2\times
  \quad\text{(worst sustained load)}, \\
  \frac{I_{\mathrm{lim}}^{\min}}{I_{\mathrm{pk}}}\bigg|_{\SI{0.5}{\ampere}}
    &= \frac{2.5}{0.74} = 3.4\times
  \quad\text{(capability point)}.
\end{align*}
This is the check that was pro-forma on the 5~V rail and is meaningful here: the
operating peak is a quarter of the limit rather than a fiftieth, but $3.4\times$ at the
full Espressif sizing point is still comfortable, and start-up adds no interaction ---
the \SI{4}{\milli\second} internal soft-start charges the output-node capacitance
(\SI{66}{\micro\farad} nominal including the module-local C502) at
\begin{align*}
  I_{\mathrm{chg}} &= \SI{66}{\micro\farad} \times
    \frac{\SI{3.3}{\volt}}{\SI{4}{\milli\second}} = \SI{54}{\milli\ampere},
\end{align*}
which even summed with the full load leaves $6\times$ under the minimum limit. A shorted
rail lands in hiccup mode, as on the 5~V rail. Match.

\subsection{Output capacitance, transient response, and rail accuracy}\label{sec:3v3-cout}

The populated $2\times\SI{22}{\micro\farad}$ (C413/C414) equals the family BOM example
and the \SIrange{22}{68}{\micro\farad} recommendation (\cref{sec:ap632xx}). The same
capacitor assumptions as \cref{sec:5v-cout} are carried --- X5R/X7R ceramics, effective
ESR \SI{2.5}{\milli\ohm} paralleled --- with DC-bias derating milder at 3.3~V:
\emph{assume} $\geq 70\,\%$ retention, i.e.\ effective buck-local
$C_{\mathrm{out}} \geq \SI{30.8}{\micro\farad}$. C501
(\SI{100}{\nano\farad}) and C502 (\SI{22}{\micro\farad}) sit on the same net at the
module: the netlist does not label intent, but their values and placement match standard
practice for a second, load-local decoupling stage at a high-transient-current device,
and they are treated as that --- C502 adds $\approx\SI{15}{\micro\farad}$ effective
bulk at the module's pins beyond the converter's own capacitors.

\paragraph{Load-release overshoot.} The datasheet's transient bound applied in
\cref{sec:5v-cout}, at the worst CCM peak and \SI{50}{\milli\volt} permitted overshoot:
\begin{align*}
  C_{\mathrm{out}} &> \frac{L\,I_{\mathrm{pk}}^2}{(\Delta V + V_{OUT})^2 - V_{OUT}^2}
    = \frac{\SI{4.7}{\micro\henry} \times 0.60^2\,\si{\square\ampere}}
           {3.35^2 - 3.30^2\ \si{\square\volt}}
    = \frac{\SI{1.67}{\micro\joule}}{\SI{0.3325}{\square\volt}}
    = \SI{5.0}{\micro\farad}
\end{align*}
(\SI{7.7}{\micro\farad} at the \SI{0.74}{\ampere} capability-point peak) --- met
$6.1\times$ over by the derated buck-local capacitance alone, $4\times$ at the
capability point. The energy argument run forward gives the hard ceiling: dumping the
full \SI{0.60}{\ampere} inductor current into \SI{30.8}{\micro\farad} raises the rail by
only $\sqrt{3.3^2 + L I^2/C} - 3.3 \approx \SI{8}{\milli\volt}$. This matters more here
than on the 5~V rail because the ESP32's absolute maximum equals its recommended
maximum, \SI{3.6}{\volt}, with no headroom above it (\cref{sec:esp32}): even so, worst
regulation maximum plus worst overshoot, $3.33 + 0.008 \approx \SI{3.34}{\volt}$,
leaves \SI{260}{\milli\volt} to the ceiling. Match.

\paragraph{Load-step droop (Wi-Fi TX burst onset).} The bounding step is
idle-to-TX, $\Delta I = \SI{0.355}{\ampere}$. Two mechanisms stack. The inductor's
physical slew limit, worst at minimum input,
\begin{align*}
  \frac{dI_L}{dt} &= \frac{V_N - V_{OUT}}{L} = \frac{8.1-3.3}{\SI{4.7}{\micro\henry}}
    = \SI{1.02}{\ampere\per\micro\second}
  &&\Rightarrow&
  t_{\mathrm{slew}} &= \frac{0.355}{1.02} = \SI{0.35}{\micro\second},
\end{align*}
costs charge $\tfrac{1}{2}\,\Delta I\, t_{\mathrm{slew}} = \SI{62}{\nano\coulomb}$,
i.e.\ \SI{2}{\milli\volt} on \SI{30.8}{\micro\farad} --- negligible; the
\SI{4.7}{\micro\henry} reuse costs nothing measurable here. The dominant term is
control-loop and PFM$\to$PWM transition latency, for which the datasheet publishes no
bandwidth. \emph{Assumption:} the loop commands full inductor current within
\SI{5}{\micro\second} ($\approx$5~switching periods), a conservative figure for an
internally compensated \SI{1.1}{\mega\hertz} peak-current-mode buck. Then
\begin{align*}
  \Delta V_{\mathrm{droop}} &\approx \frac{\Delta I \cdot t_{\mathrm{resp}}}{C_{\mathrm{out}}}
      + \Delta I \cdot \mathrm{ESR}
    = \frac{0.355 \times \SI{5}{\micro\second}}{\SI{30.8}{\micro\farad}} + 0.355\times\SI{2.5}{\milli\ohm}
    \approx 58 + 1 \approx \SI{59}{\milli\volt}.
\end{align*}
Stacked on the worst regulation minimum: $3.27 - 0.059 = \SI{3.21}{\volt}$, a
\SI{210}{\milli\volt} margin over the ESP32's \SI{3.0}{\volt} floor. The conclusion is
robust to the assumption: even tripling the assumed latency to \SI{15}{\micro\second}
(\SI{173}{\milli\volt} droop) keeps the rail at \SI{3.10}{\volt}, still in
specification, and C502's placement at the module supplies the first microseconds
locally. Match.

\paragraph{Ripple and accuracy stack.} CCM ripple at full load,
$\Delta I_L/(8 f_{SW} C) + \Delta I_L\,\mathrm{ESR} \approx 1.8 + 1.2 \approx
\SI{3}{\milli\volt}$; PFM ripple at idle, bounded by one
\SI{450}{\milli\ampere}-peak pulse's charge (method of \cref{sec:5v-cout};
$Q = \SI{243}{\nano\coulomb}$ at the \SI{8.1}{\volt} corner) at
$Q/C + I_{\mathrm{pk}}\,\mathrm{ESR} \approx 7.9 + 1.1 \approx \SI{9}{\milli\volt}$
pk-pk. The full stack against the load's window:
\begin{align*}
  V_{\mathrm{rail}} &\in [\,3.27 - 0.009,\; 3.33 + 0.009\,]
    = [\,\SI{3.26}{\volt},\ \SI{3.34}{\volt}\,]
  \quad\text{vs.}\quad \text{VDD33} \in [\,\SI{3.0}{\volt},\ \SI{3.6}{\volt}\,],
\end{align*}
$\approx\SI{260}{\milli\volt}$ of margin on both sides (\cref{sec:esp32}), and the
transient droop above consumes under a third of the low-side margin. Match.

\subsection{Input capacitance}\label{sec:3v3-cin}

C411 (\SI{10}{\micro\farad} ceramic) on the OR node equals the family BOM example
against the ``greater than \SI{10}{\micro\farad}'' prose (\cref{sec:ap632xx}), the same
precision-not-flag status as C401. The buck's input RMS ripple current at the worst
point ($D$ nearest 0.5, i.e.\ $V_N = \SI{8.1}{\volt}$):
\begin{align*}
  I_{\mathrm{RMS}} &\approx I_{\mathrm{out}}\sqrt{D(1-D)}
    = 0.359\sqrt{0.407\times0.593} = \SI{0.18}{\ampere},
\end{align*}
equal to the half-of-max-load guidance figure and trivially inside any ceramic
\SI{10}{\micro\farad}'s ampere-class rating. Input ripple voltage,
$I_{\mathrm{out}}D(1-D)/(f_{SW}C) \approx \SI{8}{\milli\volt}$ nominal
($\approx\SI{13}{\milli\volt}$ after DC-bias derating at the 12.8~V corner), is
absorbed at the node --- and, usefully, the buck's pulsed input current circulates
between C411 and the IC, so D411 carries the smoothed average
(\SI{0.24}{\ampere} worst) rather than the switching pulses. As with C401, the
schematic does not record the voltage rating: on a node reaching \SI{12.8}{\volt} it
must be $\geq\SI{16}{\volt}$, \SI{25}{\volt} preferred for derating headroom ---
confirm at BOM finalization. C412 (\SI{100}{\nano\farad}, SW-to-BST) matches the
recommended bootstrap value. Match throughout.

\subsection{Diode-OR input path and UART-powered operation}\label{sec:3v3-diodeor}

The OR node (C411, U411 IN/EN) is fed by two SOD-123 Schottky diodes: D411, anode on
$V_{\mathrm{BUS}}$, and D412, anode on \texttt{VCC\_UART} --- pin~3 of J501, the 6-pin
UART header (1/2 GND, 3 VCC, 4 RXD0, 5 TXD0, 6 unconnected). The netlist confirms
\texttt{VCC\_UART} touches exactly these two points; it reaches nothing else on the
board. The node therefore follows whichever source is higher, less one forward drop,
with the other diode reverse-biased --- a protective OR, not a sum. Both directions are
worked through below; common USB-UART adapters (FTDI/CP2102-class) ship configured for
either \SI{3.3}{\volt} or \SI{5}{\volt} on their VCC pin, often jumper-selectable, so
both sub-cases of the UART-powered scenario are treated.

\paragraph{Normal operation ($V_{\mathrm{BUS}}$ present).} The node sits at
\SIrange{8.1}{12.8}{\volt}, so D412's cathode is \SIrange{3}{13}{\volt} above any
adapter VCC (\SIrange{0}{5}{\volt}): D412 is reverse-biased and no forward path exists
from this board's power onto an attached adapter's VCC pin. The only residual coupling
is D412's reverse leakage --- \si{\micro\ampere}-class under the stated Schottky
assumption --- sourced \emph{into} the adapter's own regulated output, which is
inconsequential. Plainly: the board cannot back-power a UART adapter's VCC pin. (For
completeness on the same header: the activity-LED feeds R502/R503 can source
$\lesssim\SI{1}{\milli\ampere}$, series-limited by \SI{1.5}{\kilo\ohm}, into an
\emph{unpowered} adapter's data pins through their ESD structures --- a negligible
stress, noted only so the ``no backfeed'' statement is scoped to the VCC pin
precisely.)

\paragraph{UART-only operation, adapter set to 3.3~V.} With
$V_{\mathrm{BUS}}$ absent, D411 is reverse-biased (blocking the node from back-feeding
the bus; its \si{\micro\ampere} leakage cannot lift the bus net against the flyback's
UVLO divider) and D412 conducts. The load is the download/console-mode estimate of
\cref{sec:3v3-load}, $\approx\SI{50}{\milli\ampere}$; in any pass-through mode the
buck's input current equals its output current, so $V_F(\mathrm{D412}) \approx
\SI{0.3}{\volt}$ and
\begin{align*}
  V_N &= 3.30 - 0.30 = \SI{3.0}{\volt}.
\end{align*}
This fails before regulation is even in question: the AP63203's VIN UVLO thresholds are
\SI{3.50}{\volt} rising, $\approx\SI{3.06}{\volt}$ falling (\cref{sec:ap632xx}), so at
\SI{3.0}{\volt} the converter never enables and the +3V3 rail stays at zero. No
tolerance rescues it --- even the friendliest corner (adapter at $+5\,\%$,
\SI{3.47}{\volt}; $V_F$ at its \SI{0.25}{\volt} floor) reaches only
\SI{3.22}{\volt}, still below the rising threshold, and the recommended operating
minimum is \SI{3.8}{\volt} regardless. The counterfactual is also checked so the
conclusion does not rest on the UVLO figure alone: \emph{were} the part running, it
would be in the qualitative LDO mode of \cref{sec:ap632xx} (no numeric dropout spec
exists), with the HS FET nearly always on and the output estimated through the
resistive path:
\begin{align*}
  V_{OUT} &\approx V_N - I\,(R_{DS(\mathrm{on}),HS} + \mathrm{DCR})
    = 3.00 - 0.05\,(0.125 + 0.050) = \SI{2.99}{\volt}
\end{align*}
at \SI{50}{\milli\ampere} (\SI{2.98}{\volt} at \SI{100}{\milli\ampere}) --- below the
ESP32's \SI{3.0}{\volt} recommended minimum by \SIrange{10}{20}{\milli\volt}
(negative margin) and \SI{280}{\milli\volt} below the regulation minimum. The
\SI{3.3}{\volt} adapter configuration therefore fails both ways: the realistic outcome
is a dead rail, and even the most favorable counterfactual is out of specification.
Flagged: incompatible.

\paragraph{UART-only operation, adapter set to 5~V.} At the
\SI{50}{\milli\ampere} programming load, $I_{IN} \approx \SI{40}{\milli\ampere}$
($\eta\approx0.88$ at this light load and low ratio), $V_F \approx \SI{0.3}{\volt}$:
\begin{align*}
  V_N &= 5.0 - 0.3 = \SI{4.7}{\volt},
\end{align*}
above the \SI{3.50}{\volt} rising UVLO and the \SI{3.8}{\volt} recommended minimum,
with $D = 3.3/4.7 = 0.70$ --- deep inside the part's near-100\,\% duty capability, so
the buck regulates normally at \SIrange{3.27}{3.33}{\volt}, identical to bus-powered
operation. The margin survives every stacked tolerance: a USB-derived
``\SI{5}{\volt}'' sagging to \SI{4.5}{\volt} still yields $V_N \approx
\SI{4.1}{\volt} > \SI{3.8}{\volt}$, and even firmware that lights the radio on UART
power ($I_{IN} \approx \SI{0.30}{\ampere}$, $V_F \approx \SI{0.4}{\volt}$, $V_N
\approx \SI{4.6}{\volt}$) stays regulated --- the limit in that corner is the
adapter's own current capability, not this converter. The \SI{5}{\volt} setting is
therefore not merely the safer choice; it is the \emph{only} working one. The
intuitive pairing --- ``3.3~V board, 3.3~V adapter'' --- is exactly backwards here,
because the rail is regulated on-board and the regulator needs UVLO clearance and
headroom above its own output.

\paragraph{Conclusion and user-guide requirement.} This board's UART header requires
the adapter's VCC set or wired to \SI{5}{\volt}; a \SI{3.3}{\volt}-VCC adapter leaves
the board dead when USB-C power is absent. This is an operational requirement for the
hardware user guide's flashing/console procedure, stated there as: \emph{use an
adapter supplying \SI{5}{\volt} on VCC with \SI{3.3}{\volt} logic levels on TXD/RXD}
--- the second clause because the data pins connect directly to ESP32 GPIOs whose
recommended input maximum is $V_{DD}+\SI{0.3}{\volt}$ (\cref{sec:esp32}); many
FTDI-style cables offer precisely this 5~V-VCC/3.3~V-I/O combination. Attaching or
detaching the adapter while USB-C power is present is safe in either order: the OR
simply follows the higher source, which is always $V_{\mathrm{BUS}}$.

\paragraph{Component ratings on the OR node.} Across all scenarios the node spans
\SIrange{0}{12.8}{\volt}: U411's IN/EN absolute maximum of \SI{35}{\volt}
(\cref{sec:ap632xx}) is met $2.7\times$ over, and C411's required $\geq\SI{16}{\volt}$
rating is the confirm-at-BOM item of \cref{sec:3v3-cin}. D411 sees at most
$\approx\SI{5}{\volt}$ reverse (5~V adapter, bus absent) and \SI{0.24}{\ampere}
forward average (\SI{96}{\milli\watt} --- inside a SOD-123's typical
\SIrange{400}{500}{\milli\watt} envelope, under the stated assumption of an
$I_F\geq\SI{0.5}{\ampere}$-class part). D412 is the binding case: it blocks the full
node voltage, up to \SI{12.8}{\volt}, whenever the bus is up and no adapter drives
VCC, so the part must be rated $V_R \geq \SI{16}{\volt}$ (\SI{20}{\volt}-plus is the
commodity floor for SOD-123 Schottkys, but the schematic records no part number) ---
confirm at BOM finalization alongside the inductor and capacitor ratings.

\subsection{Pull-up loads, soft start, and dissipation}\label{sec:3v3-house}

\paragraph{PGOOD pull-up (R206).} The LM5155's open-drain PGOOD pin is pulled up by
R206 (\SI{100}{\kilo\ohm}) to this rail, with ESP32 GPIO15 --- an input, contributing
only leakage --- reading the same node. The pin-level checks (sink current
\SI{33}{\micro\ampere} vs.\ the \SI{1}{\milli\ampere} limit,
$\SI{100}{\kilo\ohm} \geq \SI{10}{\kilo\ohm}$ recommended, voltage ratings) were
completed in \cref{sec:hv-house} and stand. What remains for this rail is the load
contribution and its polarity. The pin releases to high impedance when the flyback's
output is good and its VCC and UVLO conditions are met, and pulls low otherwise
(\ds{lm5155.pdf}{22}) --- i.e.\ the pull-up conducts its
$3.3/\SI{100}{\kilo\ohm} = \SI{33}{\micro\ampere}$ only during start-up and fault
windows, and essentially nothing (GPIO leakage, \si{\nano\ampere}) in normal operation.
Either way it is four decades under the design load: negligible. GPIO15 accordingly
reads high for power-good --- a coherent arrangement, since the signal's only consumer
is powered by the same rail that supplies the pull-up, so the signal is valid exactly
when there is something running to read it.

\paragraph{EN pull-up (R501).} The R501/C503 network is Espressif's reference circuit
verbatim (\cref{sec:esp32}). Steady-state current into the EN input is leakage-level;
R501 conducts measurably only while charging C503 at power-up or while the reset
button holds EN low --- \SI{330}{\micro\ampere} peak, decaying with
$\tau = \SI{10}{\milli\second}$, or held during a button press. Zero contribution to
the steady-state budget; negligible.

\paragraph{Soft start and dissipation.} Soft start is the fixed internal
\SI{4}{\milli\second} ramp; the inrush interaction was dismissed in
\cref{sec:3v3-ilim}. Converter dissipation at the conservative $\eta=0.85$:
\begin{align*}
  P_{\mathrm{diss}} &= P_{\mathrm{out,max}}\left(\frac{1}{\eta}-1\right)
    = 1.19 \times 0.176 = \SI{0.21}{\watt}, &
  \Delta T &= \theta_{JA}P_{\mathrm{diss}} = 89 \times 0.21 \approx \SI{19}{\kelvin},
\end{align*}
rising to \SI{0.29}{\watt} and $\Delta T \approx \SI{26}{\kelvin}$ at the
\SI{0.5}{\ampere} capability point --- $T_J \leq \SI{66}{\celsius}$ at a
\SI{40}{\celsius} ambient, far below the \SI{125}{\celsius} design target
(\cref{sec:ap632xx}). Wi-Fi TX at 100\,\% duty is already the sustained rating, so no
duty-cycle relief is needed to make this conclusion. Add D411's \SI{96}{\milli\watt}
and the thermal picture on this rail remains untroubled.

\subsection{Summary}\label{sec:3v3-summary}

Assumptions carried through this section: generic SOD-123 Schottky forward drops
\SIrange{0.3}{0.45}{\volt} over \SIrange{50}{500}{\milli\ampere} with
\si{\micro\ampere}-class leakage; inductor $I_{\mathrm{sat}}\geq\SI{3.0}{\ampere}$ and
DCR $\leq\SI{50}{\milli\ohm}$ (6045-class, shared with L401); output ceramics retaining
$\geq70\,\%$ at 3.3~V bias, effective ESR \SI{2.5}{\milli\ohm}; converter efficiency
0.85 (floor under the datasheet's $\approx$87--90\,\% curve); control-loop response to
a load step within \SI{5}{\micro\second} (conclusion shown robust to $3\times$ that);
LED forward voltage \SI{2.0}{\volt}; download-mode MCU draw
$\approx\SI{50}{\milli\ampere}$ (bounded at \SI{100}{\milli\ampere}). Each appears at
its point of use above.

\begin{longtable}{@{}p{0.26\textwidth}p{0.30\textwidth}p{0.22\textwidth}p{0.14\textwidth}@{}}
\toprule
Parameter & Derived requirement & Populated & Verdict \\
\midrule
\endfirsthead
\toprule
Parameter & Derived requirement & Populated & Verdict \\
\midrule
\endhead
\bottomrule
\endlastfoot
Load current & \SI{355}{\milli\ampere} MCU $+$ $\leq$\SI{4}{\milli\ampere}
  housekeeping $\approx$ \SI{0.36}{\ampere} sustained; capability
  $\geq\SI{0.5}{\ampere}$ (Espressif) & \SI{2}{\ampere}-rated AP63203 &
  Match ($4\times$) \\
Rail vs.\ MCU window & VDD33 \SIrange{3.0}{3.6}{\volt}; abs max $=$ rec max &
  \SIrange{3.26}{3.34}{\volt} incl.\ ripple & Match
  ($\approx\SI{260}{\milli\volt}$ both sides) \\
Input window (normal) & inside \SIrange{3.8}{32}{\volt} recommended & OR node
  \SIrange{8.1}{12.8}{\volt} ($V_{\mathrm{BUS}}-V_F$) & Match \\
Duty cycle & 0.258--0.407; never crosses 0.5 & --- & Verified \\
Conduction mode & CCM boundary \SIrange{189}{237}{\milli\ampere}: PFM at idle,
  CCM in TX bursts; automatic transition & --- & Verified (benign) \\
Minimum on-time & \SI{221}{\nano\second} vs.\ \SI{80}{\nano\second}; binds only
  at $V_N\approx\SI{35}{\volt}$ & --- & Match ($2.8\times$) \\
Inductance L411 & Eq.\ window \SIrange{2.2}{3.7}{\micro\henry}; Table-2 BOM
  \SI{3.9}{\micro\henry} & \SI{4.7}{\micro\henry} (5~V rail's part reused) &
  Acceptable deviation (ripple 19--24\,\%, benign) \\
Inductor $I_{\mathrm{sat}}$ & $\geq\SI{0.68}{\ampere}$ (rule at
  \SI{0.5}{\ampere} point); fault $\leq\SI{3.1}{\ampere}$ & 6045-class, rating
  unrecorded (assumed $\geq\SI{3}{\ampere}$) & Match under assumption;
  \textbf{confirm part} \\
Current-limit margin & peak \SI{0.60}{\ampere} (\SI{0.74}{\ampere} at
  capability point) vs.\ \SI{2.5}{\ampere} min & --- & Match ($\geq3.4\times$) \\
Output capacitance C413/C414 & $\geq\SI{5.0}{\micro\farad}$ transient bound;
  \SIrange{22}{68}{\micro\farad} recommended & $2\times\SI{22}{\micro\farad}$
  ($+$ C502 at module) & Match \\
TX-burst droop & $\lesssim\SI{59}{\milli\volt}$ vs.\ \SI{270}{\milli\volt}
  budget to \SI{3.0}{\volt}; robust to $3\times$ latency & --- & Match \\
Output ripple & $\leq\SI{9}{\milli\volt}$ pk-pk PFM / $\approx\SI{3}{\milli\volt}$
  CCM & --- & Derived; $\ll$ margins \\
Input capacitor C411 & $\geq\SI{10}{\micro\farad}$; RMS $\geq\SI{0.18}{\ampere}$;
  $\geq\SI{16}{\volt}$ rating & \SI{10}{\micro\farad} (= BOM example; rating
  unrecorded) & Match; \textbf{confirm rating} \\
BST capacitor C412 & \SI{100}{\nano\farad} SW-to-BST & \SI{100}{\nano\farad} &
  Match \\
Backfeed to adapter VCC & none permissible & D412 orientation blocks; leakage
  \si{\micro\ampere}-class & Verified (none) \\
UART power, \SI{3.3}{\volt} adapter & node \SI{3.0}{\volt} $<$
  \SI{3.50}{\volt} UVLO $\Rightarrow$ rail dead; counterfactual LDO output
  \SI{2.98}{\volt}, below MCU minimum & --- & \textbf{Flagged: incompatible} \\
UART power, \SI{5}{\volt} adapter & node \SIrange{4.1}{4.7}{\volt}; normal
  regulation & --- & Match; \textbf{required configuration} (user-guide item,
  with 3.3~V logic levels) \\
D412 reverse rating & $\geq\SI{16}{\volt}$ (blocks \SI{12.8}{\volt}) & generic
  Schottky, part unrecorded & Match under assumption; \textbf{confirm part} \\
D411 forward path & $I_F$ avg $\leq\SI{0.24}{\ampere}$, \SI{96}{\milli\watt} &
  generic SOD-123 Schottky & Match under assumption \\
PGOOD pull-up R206 & load \SI{33}{\micro\ampere}, not-good windows only
  (high-Z when good) & \SI{100}{\kilo\ohm}; pin checks in \cref{sec:hv-house} &
  Match; negligible \\
EN pull-up R501 & transient only (\SI{330}{\micro\ampere} pk,
  $\tau=\SI{10}{\milli\second}$) & \SI{10}{\kilo\ohm} (= Espressif reference) &
  Match; negligible \\
Indicator LEDs R411/R502/R503 & \SI{0.87}{\milli\ampere} continuous $+$
  $\leq\SI{1.6}{\milli\ampere}$ UART activity & \SI{1.5}{\kilo\ohm} each
  (D413 value field mislabeled ``RX'') & Negligible; fix label \\
Module decoupling C501/C502 & load-local second stage at a burst load &
  \SI{100}{\nano\farad} $+$ \SI{22}{\micro\farad} & Match (role inferred) \\
Soft start / inrush & \SI{54}{\milli\ampere} charging $+$ load $\ll$ limit &
  internal \SI{4}{\milli\second} & Match \\
Converter dissipation & $\leq\SI{0.29}{\watt}$, $\Delta T \leq \SI{26}{\kelvin}$
  vs.\ \SI{125}{\celsius} target & TSOT26,
  $\theta_{JA}=89$~\si{\celsius\per\watt} & Match (large margin) \\
PSRAM active-mode adder & unquantified above \SI{355}{\milli\ampere}
  (\cref{sec:esp32}) & module SKU unrecorded & Open item (absorbed by
  $4\times$ headroom) \\
\end{longtable}

Three findings distinguish this rail from its 5~V sibling. First, the inductor is a
deliberate deviation from the variant-specific reference BOM (\SI{4.7}{\micro\henry}
where Table~2 lists \SI{3.9}{\micro\henry}) --- acceptable, with every consequence
benign at this rail's loads. Second, and operationally most significant: the UART
header's diode-OR feed works \emph{only} with a \SI{5}{\volt}-VCC adapter; a
\SI{3.3}{\volt}-VCC adapter cannot lift the buck's input past its own UVLO and leaves
the board dead when USB-C power is absent. That requirement (5~V VCC, 3.3~V logic
levels) belongs in the user guide's flashing procedure as a hard prerequisite, not a
footnote. Third, the transient checks that were pro-forma on the 5~V rail are real here
but all pass with room: worst droop uses under a third of the low-side voltage margin,
and the current-limit margin at Espressif's full \SI{0.5}{\ampere} sizing point is
$3.4\times$. The confirm-at-BOM items --- inductor saturation rating, C411 voltage
rating, D412 reverse rating --- are unrecorded ratings on generic parts, the same
status as their counterparts in \cref{sec:5v-converter}, plus one cosmetic fix (D413's
value field) and one inherited open item (the PSRAM active-mode current adder).
